\subsubsection{二阶标量方程}\label{subsec:second-order-scalar-interior-regularity}

为免去许多技术细节, 只对二阶标量方程给出内部正则性结果的证明. 前面 Section~\ref{sec:general-linear-elliptic-problems} 已经陈述了高阶方程和方程组的结果.

\begin{Theorem}{内部正则性}
\label{thm:second-order-interior-regularity}
设 $L$ 为形如下式的一致椭圆二阶算子\[L(\bm{x},D)u:=\sum_{i,j=1}^n\frac{\partial}{\partial x_i}a_{ij}(\bm{x})\frac{\partial u}{\partial x_j}+\sum_{i=1}^nb_i(\bm{x})\frac{\partial u}{\partial x_i}+c(\bm{x})u,\]对应的双线性型为\[B[v,u]:=-\sum_{i,j=1}^n\int_\Omega a_{ij}(\bm{x})u_{x_j}v_{x_i}\,\mathrm{d}\bm{x}+\sum_{i=1}^n\int_\Omega b_i(\bm{x})u_{x_i}v\,\mathrm{d}\bm{x}+\int_\Omega c(\bm{x})uv\,\mathrm{d}\bm{x}.\]假设系数满足 $a_{ij}\in W^{1,\infty}(\Omega)$, $b_i,c\in L^\infty(\Omega)$, 且 $f\in L^2(\Omega)$. 设 $u\in H_0^1(\Omega)$ 为 $L(D,\bm{x})u=f$ 的 Dirichlet 问题的弱解. 则对每个 $\Omega^{\prime}\subset\subset\Omega$, 有 $u\in H^2(\Omega^{\prime})$, 且
\begin{equation}\label{eq:second-order-interior-regularity-h2-estimate}
\|u\|_{H^2(\Omega^{\prime})}\leq C(\|u\|_{H^1(\Omega)}+\|f\|_{L^2(\Omega)}).
\end{equation}

\end{Theorem}

\begingroup
\renewcommand{\endProof}{\par\nobreak\hfill$\blacksquare$\par\vspace{0.1cm}}
\begin{Proof}
从恒等式
\begin{equation}\label{eq:interior-regularity-principal-weak-identity}
\int_\Omega a_{ij}u_{x_j}v_{x_i}\,\mathrm{d}\bm{x}=\int_\Omega gv\,\mathrm{d}\bm{x}
\end{equation}
开始, 它对所有 $v\in H_0^1(\Omega)$ 成立, 其中 $g\in L^2(\Omega)$ 由
\begin{equation}\label{eq:interior-regularity-reduced-right-hand-side}
g:=b_iu_{x_i}+cu-f.
\end{equation}
给出. 现在假设 $v\in H^1(\Omega)$ 且 $\operatorname{supp}v\subset\subset\Omega$, 并取\[2|h|<\operatorname{dist}(\operatorname{supp}v,\partial\Omega).\]则对任意 $k=1,\ldots,n$, 有 $D_k^{-h}v\in H_0^1(\Omega)$, 因而可利用 \eqref{eq:interior-regularity-principal-weak-identity} 和 ``差分分部求和公式\textquotesingle\textquotesingle\ \eqref{eq:difference-quotient-integration-by-parts} 得到\[\begin{aligned}
\int_\Omega D_k^h(a_{ij}u_{x_j})v_{x_i}\,\mathrm{d}\bm{x}
&=-\int_\Omega a_{ij}u_{x_j}D_k^{-h}v_{x_i}\,\mathrm{d}\bm{x}\\
&=-\int_\Omega gD_k^{-h}v\,\mathrm{d}\bm{x}.
\end{aligned}\]再利用这一等式和差商的乘积法则 \eqref{eq:difference-quotient-shifted-product-rule}, 得到
\begin{equation}\label{eq:interior-regularity-shifted-principal-product-identity}
\int_\Omega a_{ij}(\bm{x}+h\bm{e}_k)(D_k^hu_{x_j})v_{x_i}\,\mathrm{d}\bm{x}=-\int_\Omega(D_k^ha_{ij})u_{x_j}v_{x_i}+gD_k^{-h}v\,\mathrm{d}\bm{x}.
\end{equation}
现在利用 Lemma~\ref{lem:difference-quotient-weak-derivative-bound}, 可将其估计为
\begin{equation}\label{eq:interior-regularity-shifted-principal-form-bound}
\left|\int_\Omega a_{ij}(\bm{x}+h\bm{e}_k)(D_k^hu_{x_j})v_{x_i}\,\mathrm{d}\bm{x}\right|\leq C(\|u\|_{1,2}+\|f\|_2)\|\nabla v\|_2.
\end{equation}
设 $\Omega^{\prime}\subset\subset\Omega$. 选取截断函数 $\zeta\in\mathcal{D}(\Omega)$, 具有以下性质:\begin{enumerate}\item 在 $\Omega^{\prime}$ 上 $\zeta\equiv1$.\item $|\nabla\zeta|<2/d$, 其中 $d=\operatorname{dist}(\Omega^{\prime},\partial\Omega)$.\end{enumerate}现在利用椭圆性条件 \eqref{eq:second-order-ellipticity-negative-coefficient-matrix}, 初等不等式和前面的估计 (取 $v=\zeta^2D_k^hu$), 得到\[\begin{aligned}
\theta\int_\Omega|\zeta\nabla D_k^hu|^2\,\mathrm{d}\bm{x}
&\leq-\int_\Omega\zeta^2a_{ij}(\bm{x}+h\bm{e}_k)D_k^hu_{x_i}D_k^hu_{x_j}\,\mathrm{d}\bm{x}\\
&=-\int_\Omega a_{ij}(\bm{x}+h\bm{e}_k)D_k^hu_{x_i}\left[(\zeta^2D_k^hu)_{x_j}-2D_k^hu\zeta\zeta_{x_j}\right]\,\mathrm{d}\bm{x}\\
&\leq C\Bigl[(\|u\|_{1,2}+\|f\|_2)(\|\zeta\nabla D_k^hu\|_2+2\|D_k^hu\nabla\zeta\|_2)\\
&\qquad+\sup_{i,j}\|a_{ij}\|_\infty\|\zeta\nabla D_k^hu\|_2\|D_k^hu\nabla\zeta\|_2\Bigr]\\
&\leq C(\epsilon)[(\|u\|_{1,2}+\|f\|_2+\|D_k^hu\nabla\zeta\|_2)^2+\epsilon\|\zeta\nabla D_k^hu\|_2^2]
\end{aligned}\]对任意 $\epsilon>0$ 成立.

因此, 适当选取 $\epsilon$ 并整理后, 有\[\begin{aligned}
\|D_k^h\nabla u\|_{L^2(\Omega^{\prime})}&\leq\|\zeta D_k^h\nabla u\|_{L^2(\Omega)}\\
&\leq C(\|u\|_{H^1(\Omega)}+\|f\|_{L^2(\Omega)}+\|D_k^hu\nabla\zeta\|)\\
&\leq C(\|u\|_{H^1(\Omega)}+\|f\|_{L^2(\Omega)}).
\end{aligned}\]这里利用了 Lemma~\ref{lem:difference-quotient-weak-derivative-bound} 和对 $|\nabla\zeta|$ 的逐点界.

最后, 利用关于 $\nabla u$ 的差商的这一估计和 Lemma~\ref{lem:uniform-difference-quotient-converse}, 可推出 $\nabla u\in H^1(\Omega^{\prime})$. 估计 \eqref{eq:second-order-interior-regularity-h2-estimate} 立即得到.
\end{Proof}
\endgroup
